%%%%%%%%%%%%%%%%%%%%%%%%%%%%%%%%%%%%%%%%%%%%%%%%%%%%%%%%%%%%
% 10.11 QUEUEING THEORY
% The Composition of Advanced Mathematics
%%%%%%%%%%%%%%%%%%%%%%%%%%%%%%%%%%%%%%%%%%%%%%%%%%%%%%%%%%%%

\section{Queueing Theory}
\label{sec:queueing-theory}

\prerequisite{
Probability, random variables, expectation, exponential and Poisson
distributions, stochastic processes, Markov chains, differential
equations, and basic optimization.
}

\learningobjectives{
By the end of this section, the reader should be able to:
\begin{itemize}
    \item define a queueing system;
    \item identify customers, servers, arrivals, service times, and queues;
    \item distinguish arrival processes from service processes;
    \item understand queue capacity and calling populations;
    \item recognize common queue disciplines;
    \item interpret Kendall notation;
    \item understand Poisson arrivals;
    \item understand exponential service times;
    \item explain the memoryless property;
    \item define arrival and service rates;
    \item define utilization and traffic intensity;
    \item derive steady-state probabilities for birth--death queues;
    \item analyze the \(M/M/1\) queue;
    \item compute \(L\), \(L_q\), \(W\), and \(W_q\);
    \item use Little's Law;
    \item understand stability conditions;
    \item analyze finite-capacity \(M/M/1/K\) systems conceptually;
    \item analyze \(M/M/c\) multi-server queues;
    \item understand Erlang delay formulas;
    \item distinguish Erlang B and Erlang C models;
    \item understand loss systems;
    \item recognize \(M/G/1\) queues;
    \item state the Pollaczek--Khinchine mean waiting-time formula;
    \item understand deterministic service as a special case;
    \item recognize \(G/G/1\) systems;
    \item understand variability effects on waiting time;
    \item recognize Kingman's approximation;
    \item understand finite-source queueing systems conceptually;
    \item distinguish transient and steady-state analysis;
    \item understand queueing networks conceptually;
    \item recognize Jackson networks;
    \item understand bottleneck utilization;
    \item understand priority queues conceptually;
    \item distinguish preemptive and nonpreemptive priorities;
    \item recognize scheduling and queueing connections;
    \item understand balking, reneging, and jockeying;
    \item recognize service-level measures;
    \item understand staffing tradeoffs;
    \item connect queueing theory with operations research,
          telecommunications, healthcare, transportation, manufacturing,
          and computing.
\end{itemize}
}

%%%%%%%%%%%%%%%%%%%%%%%%%%%%%%%%%%%%%%%%%%%%%%%%%%%%%%%%%%%%
\subsection{What Is a Queueing System?}
\label{subsec:what-is-queueing-system}
%%%%%%%%%%%%%%%%%%%%%%%%%%%%%%%%%%%%%%%%%%%%%%%%%%%%%%%%%%%%

A queueing system is a stochastic service system in which entities arrive,
possibly wait, receive service, and depart.

Examples include:

\[
\boxed{
\text{customers waiting for a cashier},
}
\]

\[
\boxed{
\text{aircraft waiting for a gate},
}
\]

\[
\boxed{
\text{jobs waiting for a machine},
}
\]

\[
\boxed{
\text{data packets waiting for transmission},
}
\]

and

\[
\boxed{
\text{patients waiting for treatment}.
}
\]

%%%%%%%%%%%%%%%%%%%%%%%%%%%%%%%%%%%%%%%%%%%%%%%%%%%%%%%%%%%%
\subsection{Basic Components of a Queue}
\label{subsec:basic-components-queue}
%%%%%%%%%%%%%%%%%%%%%%%%%%%%%%%%%%%%%%%%%%%%%%%%%%%%%%%%%%%%

A queueing model usually specifies:

\[
\boxed{
\text{arrival process},
}
\]

\[
\boxed{
\text{service-time distribution},
}
\]

\[
\boxed{
\text{number of servers},
}
\]

\[
\boxed{
\text{queue capacity},
}
\]

\[
\boxed{
\text{calling population},
}
\]

and

\[
\boxed{
\text{service discipline}.
}
\]

%%%%%%%%%%%%%%%%%%%%%%%%%%%%%%%%%%%%%%%%%%%%%%%%%%%%%%%%%%%%
\subsection{Customers and Jobs}
\label{subsec:queue-customers-jobs}
%%%%%%%%%%%%%%%%%%%%%%%%%%%%%%%%%%%%%%%%%%%%%%%%%%%%%%%%%%%%

The word \emph{customer} is used generically.

A customer may represent:

\[
\boxed{
\text{a person},
}
\]

\[
\boxed{
\text{a vehicle},
}
\]

\[
\boxed{
\text{a machine-repair request},
}
\]

\[
\boxed{
\text{a computer task},
}
\]

or

\[
\boxed{
\text{a message packet}.
}
\]

%%%%%%%%%%%%%%%%%%%%%%%%%%%%%%%%%%%%%%%%%%%%%%%%%%%%%%%%%%%%
\subsection{Arrival Process}
\label{subsec:queue-arrival-process}
%%%%%%%%%%%%%%%%%%%%%%%%%%%%%%%%%%%%%%%%%%%%%%%%%%%%%%%%%%%%

The arrival process describes when customers enter the system.

If arrivals occur at times

\[
T_1,T_2,\ldots,
\]

then the interarrival times are

\[
\boxed{
A_n
=
T_n-T_{n-1}.
}
\]

An arrival model specifies the distribution or dependence structure of the
\(A_n\).

%%%%%%%%%%%%%%%%%%%%%%%%%%%%%%%%%%%%%%%%%%%%%%%%%%%%%%%%%%%%
\subsection{Arrival Rate}
\label{subsec:queue-arrival-rate}
%%%%%%%%%%%%%%%%%%%%%%%%%%%%%%%%%%%%%%%%%%%%%%%%%%%%%%%%%%%%

The average arrival rate is commonly denoted

\[
\boxed{
\lambda.
}
\]

Its units might be:

\[
\boxed{
\text{customers per hour},
}
\]

or

\[
\boxed{
\text{jobs per second}.
}
\]

If mean interarrival time is

\[
\mathbb{E}[A],
\]

then under suitable renewal assumptions,

\[
\boxed{
\lambda
=
\frac{1}{\mathbb{E}[A]}.
}
\]

%%%%%%%%%%%%%%%%%%%%%%%%%%%%%%%%%%%%%%%%%%%%%%%%%%%%%%%%%%%%
\subsection{Service Process}
\label{subsec:queue-service-process}
%%%%%%%%%%%%%%%%%%%%%%%%%%%%%%%%%%%%%%%%%%%%%%%%%%%%%%%%%%%%

Let

\[
S_n
\]

be the service time required by customer \(n\).

The service-time distribution describes variability in processing
duration.

The average service rate of one server is often denoted

\[
\boxed{
\mu.
}
\]

If

\[
\mathbb{E}[S]
\]

is the mean service time, then

\[
\boxed{
\mu
=
\frac{1}{\mathbb{E}[S]}.
}
\]

%%%%%%%%%%%%%%%%%%%%%%%%%%%%%%%%%%%%%%%%%%%%%%%%%%%%%%%%%%%%
\subsection{Number of Servers}
\label{subsec:queue-number-servers}
%%%%%%%%%%%%%%%%%%%%%%%%%%%%%%%%%%%%%%%%%%%%%%%%%%%%%%%%%%%%

A queue may have one or several parallel servers.

Examples include:

\[
\boxed{
c=1
}
\]

for one cashier,

and

\[
\boxed{
c=5
}
\]

for five identical service counters.

The number of servers strongly affects waiting time and utilization.

%%%%%%%%%%%%%%%%%%%%%%%%%%%%%%%%%%%%%%%%%%%%%%%%%%%%%%%%%%%%
\subsection{Queue Capacity}
\label{subsec:queue-capacity}
%%%%%%%%%%%%%%%%%%%%%%%%%%%%%%%%%%%%%%%%%%%%%%%%%%%%%%%%%%%%

Queue capacity specifies the maximum number of customers allowed in the
system or waiting area.

An infinite-capacity model assumes

\[
\boxed{
K=\infty.
}
\]

A finite-capacity model may allow only

\[
\boxed{
K
}
\]

customers in the entire system.

Additional arrivals may then be blocked or lost.

%%%%%%%%%%%%%%%%%%%%%%%%%%%%%%%%%%%%%%%%%%%%%%%%%%%%%%%%%%%%
\subsection{Calling Population}
\label{subsec:calling-population}
%%%%%%%%%%%%%%%%%%%%%%%%%%%%%%%%%%%%%%%%%%%%%%%%%%%%%%%%%%%%

The calling population is the population from which arrivals originate.

An infinite-source approximation assumes that the arrival rate does not
depend substantially on the number already in the queue.

A finite-source model is useful when only a limited number of entities can
generate service requests.

%%%%%%%%%%%%%%%%%%%%%%%%%%%%%%%%%%%%%%%%%%%%%%%%%%%%%%%%%%%%
\subsection{Queue Discipline}
\label{subsec:queue-discipline}
%%%%%%%%%%%%%%%%%%%%%%%%%%%%%%%%%%%%%%%%%%%%%%%%%%%%%%%%%%%%

The queue discipline specifies how waiting customers are selected for
service.

Common disciplines include:

\[
\boxed{
\text{FCFS},
}
\]

\[
\boxed{
\text{LCFS},
}
\]

\[
\boxed{
\text{priority service},
}
\]

and

\[
\boxed{
\text{processor sharing}.
}
\]

%%%%%%%%%%%%%%%%%%%%%%%%%%%%%%%%%%%%%%%%%%%%%%%%%%%%%%%%%%%%
\subsection{First-Come, First-Served}
\label{subsec:fcfs}
%%%%%%%%%%%%%%%%%%%%%%%%%%%%%%%%%%%%%%%%%%%%%%%%%%%%%%%%%%%%

First-come, first-served, abbreviated FCFS, serves customers in order of
arrival.

It is also commonly called FIFO:

\[
\boxed{
\text{first in, first out}.
}
\]

Many standard queue formulas assume FCFS unless otherwise stated.

%%%%%%%%%%%%%%%%%%%%%%%%%%%%%%%%%%%%%%%%%%%%%%%%%%%%%%%%%%%%
\subsection{Last-Come, First-Served}
\label{subsec:lcfs}
%%%%%%%%%%%%%%%%%%%%%%%%%%%%%%%%%%%%%%%%%%%%%%%%%%%%%%%%%%%%

Last-come, first-served, abbreviated LCFS, gives service priority to the
most recent arrival.

It is also called LIFO:

\[
\boxed{
\text{last in, first out}.
}
\]

The effect of LCFS depends on whether service is preemptive.

%%%%%%%%%%%%%%%%%%%%%%%%%%%%%%%%%%%%%%%%%%%%%%%%%%%%%%%%%%%%
\subsection{Kendall Notation}
\label{subsec:kendall-notation}
%%%%%%%%%%%%%%%%%%%%%%%%%%%%%%%%%%%%%%%%%%%%%%%%%%%%%%%%%%%%

Queueing systems are often summarized using Kendall notation:

\[
\boxed{
A/S/c.
}
\]

Here:

\[
A
=
\text{arrival-time distribution code},
\]

\[
S
=
\text{service-time distribution code},
\]

and

\[
c
=
\text{number of servers}.
\]

Extended forms may also include capacity, population, and discipline.

%%%%%%%%%%%%%%%%%%%%%%%%%%%%%%%%%%%%%%%%%%%%%%%%%%%%%%%%%%%%
\subsection{Common Kendall Symbols}
\label{subsec:kendall-symbols}
%%%%%%%%%%%%%%%%%%%%%%%%%%%%%%%%%%%%%%%%%%%%%%%%%%%%%%%%%%%%

Common symbols include:

\[
\boxed{
M
=
\text{Markovian or exponential},
}
\]

\[
\boxed{
D
=
\text{deterministic},
}
\]

\[
\boxed{
G
=
\text{general distribution}.
}
\]

Thus:

\[
\boxed{
M/M/1
}
\]

means Poisson arrivals, exponential service, and one server.

%%%%%%%%%%%%%%%%%%%%%%%%%%%%%%%%%%%%%%%%%%%%%%%%%%%%%%%%%%%%
\subsection{Poisson Arrivals}
\label{subsec:poisson-arrivals}
%%%%%%%%%%%%%%%%%%%%%%%%%%%%%%%%%%%%%%%%%%%%%%%%%%%%%%%%%%%%

A Poisson arrival process with rate \(\lambda\) satisfies

\[
\boxed{
N(t)
\sim
\operatorname{Poisson}(\lambda t).
}
\]

Hence,

\[
\boxed{
\mathbb{P}(N(t)=n)
=
e^{-\lambda t}
\frac{
(\lambda t)^n
}{
n!
}.
}
\]

The mean number of arrivals during time \(t\) is

\[
\boxed{
\mathbb{E}[N(t)]
=
\lambda t.
}
\]

%%%%%%%%%%%%%%%%%%%%%%%%%%%%%%%%%%%%%%%%%%%%%%%%%%%%%%%%%%%%
\subsection{Exponential Interarrival Times}
\label{subsec:exponential-interarrival-times}
%%%%%%%%%%%%%%%%%%%%%%%%%%%%%%%%%%%%%%%%%%%%%%%%%%%%%%%%%%%%

For a Poisson process, interarrival times are independent exponential
random variables:

\[
\boxed{
A
\sim
\operatorname{Exp}(\lambda).
}
\]

Their density is

\[
\boxed{
f_A(t)
=
\lambda e^{-\lambda t},
\qquad
t\geq0.
}
\]

The mean is

\[
\boxed{
\mathbb{E}[A]
=
\frac1\lambda.
}
\]

%%%%%%%%%%%%%%%%%%%%%%%%%%%%%%%%%%%%%%%%%%%%%%%%%%%%%%%%%%%%
\subsection{Exponential Service Times}
\label{subsec:exponential-service-times}
%%%%%%%%%%%%%%%%%%%%%%%%%%%%%%%%%%%%%%%%%%%%%%%%%%%%%%%%%%%%

In an \(M/M/1\) queue,

\[
\boxed{
S
\sim
\operatorname{Exp}(\mu).
}
\]

Thus,

\[
\boxed{
\mathbb{E}[S]
=
\frac1\mu.
}
\]

The probability that service exceeds time \(t\) is

\[
\boxed{
\mathbb{P}(S>t)
=
e^{-\mu t}.
}
\]

%%%%%%%%%%%%%%%%%%%%%%%%%%%%%%%%%%%%%%%%%%%%%%%%%%%%%%%%%%%%
\subsection{Memoryless Property}
\label{subsec:queue-memoryless-property}
%%%%%%%%%%%%%%%%%%%%%%%%%%%%%%%%%%%%%%%%%%%%%%%%%%%%%%%%%%%%

The exponential distribution satisfies

\[
\boxed{
\mathbb{P}(S>s+t\mid S>s)
=
\mathbb{P}(S>t).
}
\]

Thus elapsed service time does not change the conditional distribution of
remaining service.

This memoryless property allows the number in the system to form a
continuous-time Markov chain in \(M/M/c\) models.

%%%%%%%%%%%%%%%%%%%%%%%%%%%%%%%%%%%%%%%%%%%%%%%%%%%%%%%%%%%%
\subsection{Traffic Intensity}
\label{subsec:traffic-intensity}
%%%%%%%%%%%%%%%%%%%%%%%%%%%%%%%%%%%%%%%%%%%%%%%%%%%%%%%%%%%%

For an \(M/M/1\) queue, define

\[
\boxed{
\rho
=
\frac{\lambda}{\mu}.
}
\]

This is the traffic intensity or server utilization.

For a stable \(M/M/1\) queue, we require

\[
\boxed{
\rho<1.
}
\]

%%%%%%%%%%%%%%%%%%%%%%%%%%%%%%%%%%%%%%%%%%%%%%%%%%%%%%%%%%%%
\subsection{Interpretation of Utilization}
\label{subsec:queue-utilization}
%%%%%%%%%%%%%%%%%%%%%%%%%%%%%%%%%%%%%%%%%%%%%%%%%%%%%%%%%%%%

For an \(M/M/1\) queue,

\[
\rho
\]

is the long-run fraction of time the server is busy.

Thus,

\[
\boxed{
\rho=0.8
}
\]

means the server is busy approximately \(80\%\) of the time in steady
state.

High utilization generally causes rapidly increasing waiting times.

%%%%%%%%%%%%%%%%%%%%%%%%%%%%%%%%%%%%%%%%%%%%%%%%%%%%%%%%%%%%
\subsection{Why Stability Matters}
\label{subsec:queue-stability}
%%%%%%%%%%%%%%%%%%%%%%%%%%%%%%%%%%%%%%%%%%%%%%%%%%%%%%%%%%%%

If

\[
\lambda\geq\mu
\]

in an infinite-capacity \(M/M/1\) queue, arrivals occur at least as fast as
the server can process them on average.

Then no proper steady-state queue-length distribution exists.

The queue tends to grow without bound.

Thus,

\[
\boxed{
\lambda<\mu
}
\]

is the stability condition.

%%%%%%%%%%%%%%%%%%%%%%%%%%%%%%%%%%%%%%%%%%%%%%%%%%%%%%%%%%%%
\subsection{Birth--Death Process}
\label{subsec:birth-death-queue}
%%%%%%%%%%%%%%%%%%%%%%%%%%%%%%%%%%%%%%%%%%%%%%%%%%%%%%%%%%%%

Let

\[
N(t)
=
\text{number of customers in the system at time }t.
\]

For an \(M/M/1\) queue:

\[
n
\to
n+1
\]

occurs at rate

\[
\boxed{
\lambda,
}
\]

and

\[
n
\to
n-1
\]

occurs at rate

\[
\boxed{
\mu
}
\]

when \(n\geq1\).

Thus \(N(t)\) is a birth--death process.

%%%%%%%%%%%%%%%%%%%%%%%%%%%%%%%%%%%%%%%%%%%%%%%%%%%%%%%%%%%%
\subsection{Steady-State Balance for \(M/M/1\)}
\label{subsec:mm1-balance}
%%%%%%%%%%%%%%%%%%%%%%%%%%%%%%%%%%%%%%%%%%%%%%%%%%%%%%%%%%%%

Let

\[
p_n
=
\mathbb{P}(N=n)
\]

in steady state.

Local balance gives

\[
\boxed{
\lambda p_n
=
\mu p_{n+1}.
}
\]

Therefore,

\[
p_{n+1}
=
\frac{\lambda}{\mu}
p_n
=
\rho p_n.
\]

Hence,

\[
\boxed{
p_n
=
\rho^n p_0.
}
\]

%%%%%%%%%%%%%%%%%%%%%%%%%%%%%%%%%%%%%%%%%%%%%%%%%%%%%%%%%%%%
\subsection{Normalization for \(M/M/1\)}
\label{subsec:mm1-normalization}
%%%%%%%%%%%%%%%%%%%%%%%%%%%%%%%%%%%%%%%%%%%%%%%%%%%%%%%%%%%%

Since probabilities sum to \(1\),

\[
\sum_{n=0}^{\infty}
p_n
=
1.
\]

Therefore,

\[
p_0
\sum_{n=0}^{\infty}
\rho^n
=
1.
\]

For

\[
\rho<1,
\]

\[
\sum_{n=0}^{\infty}
\rho^n
=
\frac1{1-\rho}.
\]

Thus,

\[
\boxed{
p_0
=
1-\rho.
}
\]

%%%%%%%%%%%%%%%%%%%%%%%%%%%%%%%%%%%%%%%%%%%%%%%%%%%%%%%%%%%%
\subsection{\(M/M/1\) Steady-State Distribution}
\label{subsec:mm1-steady-state-distribution}
%%%%%%%%%%%%%%%%%%%%%%%%%%%%%%%%%%%%%%%%%%%%%%%%%%%%%%%%%%%%

The stationary distribution is

\[
\boxed{
p_n
=
(1-\rho)\rho^n,
\qquad
n=0,1,2,\ldots
}
\]

for

\[
\boxed{
0\leq\rho<1.
}
\]

Thus the number of customers in the system has a geometric distribution.

%%%%%%%%%%%%%%%%%%%%%%%%%%%%%%%%%%%%%%%%%%%%%%%%%%%%%%%%%%%%
\subsection{Probability the System Is Empty}
\label{subsec:mm1-empty}
%%%%%%%%%%%%%%%%%%%%%%%%%%%%%%%%%%%%%%%%%%%%%%%%%%%%%%%%%%%%

For \(M/M/1\),

\[
\boxed{
P_0
=
1-\rho.
}
\]

Therefore the probability the server is busy is

\[
\boxed{
1-P_0
=
\rho.
}
\]

This agrees with the utilization interpretation.

%%%%%%%%%%%%%%%%%%%%%%%%%%%%%%%%%%%%%%%%%%%%%%%%%%%%%%%%%%%%
\subsection{Probability of At Least \(n\) Customers}
\label{subsec:mm1-tail}
%%%%%%%%%%%%%%%%%%%%%%%%%%%%%%%%%%%%%%%%%%%%%%%%%%%%%%%%%%%%

For the geometric stationary distribution,

\[
\mathbb{P}(N\geq n)
=
\sum_{k=n}^{\infty}
(1-\rho)\rho^k.
\]

Therefore,

\[
\boxed{
\mathbb{P}(N\geq n)
=
\rho^n.
}
\]

This gives a simple queue-length tail probability.

%%%%%%%%%%%%%%%%%%%%%%%%%%%%%%%%%%%%%%%%%%%%%%%%%%%%%%%%%%%%
\subsection{Expected Number in an \(M/M/1\) System}
\label{subsec:mm1-L}
%%%%%%%%%%%%%%%%%%%%%%%%%%%%%%%%%%%%%%%%%%%%%%%%%%%%%%%%%%%%

Let

\[
L
=
\mathbb{E}[N].
\]

Using the geometric distribution,

\[
L
=
\sum_{n=0}^{\infty}
n(1-\rho)\rho^n.
\]

The result is

\[
\boxed{
L
=
\frac{\rho}{1-\rho}.
}
\]

Since

\[
\rho
=
\frac{\lambda}{\mu},
\]

this can also be written

\[
\boxed{
L
=
\frac{\lambda}{\mu-\lambda}.
}
\]

%%%%%%%%%%%%%%%%%%%%%%%%%%%%%%%%%%%%%%%%%%%%%%%%%%%%%%%%%%%%
\subsection{Expected Number Waiting}
\label{subsec:mm1-Lq}
%%%%%%%%%%%%%%%%%%%%%%%%%%%%%%%%%%%%%%%%%%%%%%%%%%%%%%%%%%%%

The expected number in service is the probability the server is busy:

\[
\boxed{
\rho.
}
\]

Therefore,

\[
L_q
=
L-\rho.
\]

Hence,

\[
\boxed{
L_q
=
\frac{\rho^2}{1-\rho}.
}
\]

Equivalently,

\[
\boxed{
L_q
=
\frac{\lambda^2}{
\mu(\mu-\lambda)
}.
}
\]

%%%%%%%%%%%%%%%%%%%%%%%%%%%%%%%%%%%%%%%%%%%%%%%%%%%%%%%%%%%%
\subsection{Little's Law}
\label{subsec:littles-law}
%%%%%%%%%%%%%%%%%%%%%%%%%%%%%%%%%%%%%%%%%%%%%%%%%%%%%%%%%%%%

\begin{theorem}[Little's Law]
Under broad stability conditions,

\[
\boxed{
L
=
\lambda W,
}
\]

where

\[
L
=
\text{average number in the system},
\]

\[
\lambda
=
\text{effective arrival rate},
\]

and

\[
W
=
\text{average time in the system}.
\]
\end{theorem}

For the waiting line alone,

\[
\boxed{
L_q
=
\lambda W_q.
}
\]

%%%%%%%%%%%%%%%%%%%%%%%%%%%%%%%%%%%%%%%%%%%%%%%%%%%%%%%%%%%%
\subsection{Why Little's Law Is Important}
\label{subsec:why-littles-law-important}
%%%%%%%%%%%%%%%%%%%%%%%%%%%%%%%%%%%%%%%%%%%%%%%%%%%%%%%%%%%%

Little's Law does not require exponential service or Poisson arrivals.

It applies under broad conditions whenever consistent long-run averages
exist.

It connects:

\[
\boxed{
\text{inventory in system}
}
\]

with

\[
\boxed{
\text{throughput}
}
\]

and

\[
\boxed{
\text{time in system}.
}
\]

%%%%%%%%%%%%%%%%%%%%%%%%%%%%%%%%%%%%%%%%%%%%%%%%%%%%%%%%%%%%
\subsection{Mean Time in an \(M/M/1\) System}
\label{subsec:mm1-W}
%%%%%%%%%%%%%%%%%%%%%%%%%%%%%%%%%%%%%%%%%%%%%%%%%%%%%%%%%%%%

Using

\[
L=\lambda W,
\]

we obtain

\[
W
=
\frac{L}{\lambda}.
\]

Since

\[
L
=
\frac{\lambda}{\mu-\lambda},
\]

\[
\boxed{
W
=
\frac1{\mu-\lambda}.
}
\]

This includes both waiting and service time.

%%%%%%%%%%%%%%%%%%%%%%%%%%%%%%%%%%%%%%%%%%%%%%%%%%%%%%%%%%%%
\subsection{Mean Waiting Time in Queue}
\label{subsec:mm1-Wq}
%%%%%%%%%%%%%%%%%%%%%%%%%%%%%%%%%%%%%%%%%%%%%%%%%%%%%%%%%%%%

Using

\[
L_q
=
\lambda W_q,
\]

we obtain

\[
\boxed{
W_q
=
\frac{
\lambda
}{
\mu(\mu-\lambda)
}.
}
\]

Equivalently,

\[
\boxed{
W_q
=
\frac{
\rho
}{
\mu-\lambda
}.
}
\]

Since

\[
W
=
W_q+\frac1\mu,
\]

the formulas are consistent.

%%%%%%%%%%%%%%%%%%%%%%%%%%%%%%%%%%%%%%%%%%%%%%%%%%%%%%%%%%%%
\subsection{\(M/M/1\) Formula Summary}
\label{subsec:mm1-formula-summary}
%%%%%%%%%%%%%%%%%%%%%%%%%%%%%%%%%%%%%%%%%%%%%%%%%%%%%%%%%%%%

For

\[
\boxed{
\rho
=
\frac{\lambda}{\mu}
<
1,
}
\]

the \(M/M/1\) steady-state formulas are

\[
\boxed{
P_n
=
(1-\rho)\rho^n,
}
\]

\[
\boxed{
L
=
\frac{\rho}{1-\rho},
}
\]

\[
\boxed{
L_q
=
\frac{\rho^2}{1-\rho},
}
\]

\[
\boxed{
W
=
\frac1{\mu-\lambda},
}
\]

and

\[
\boxed{
W_q
=
\frac{\lambda}{
\mu(\mu-\lambda)
}.
}
\]

%%%%%%%%%%%%%%%%%%%%%%%%%%%%%%%%%%%%%%%%%%%%%%%%%%%%%%%%%%%%
\subsection{Worked Example: \(M/M/1\)}
\label{subsec:worked-mm1}
%%%%%%%%%%%%%%%%%%%%%%%%%%%%%%%%%%%%%%%%%%%%%%%%%%%%%%%%%%%%

\begin{workedexamplebox}{Single Service Counter}

Suppose customers arrive at rate

\[
\lambda=6
\]

per hour.

The server completes service at rate

\[
\mu=10
\]

per hour.

The utilization is

\[
\rho
=
\frac6{10}
=
0.6.
\]

The expected number in the system is

\[
L
=
\frac{0.6}{1-0.6}
=
1.5.
\]

The expected number waiting is

\[
L_q
=
\frac{0.6^2}{1-0.6}
=
0.9.
\]

The expected time in the system is

\[
W
=
\frac1{10-6}
=
\frac14
\text{ hour}.
\]

Thus,

\[
W=15
\text{ minutes}.
\]

The expected waiting time is

\[
W_q
=
\frac6{
10(10-6)
}
=
0.15
\text{ hour}.
\]

Thus,

\[
W_q=9
\text{ minutes}.
\]

\begin{answerbox}
\[
\boxed{
\rho=0.6,
\quad
L=1.5,
\quad
L_q=0.9,
}
\]

\[
\boxed{
W=15\text{ minutes},
\quad
W_q=9\text{ minutes}.
}
\]
\end{answerbox}

\end{workedexamplebox}

%%%%%%%%%%%%%%%%%%%%%%%%%%%%%%%%%%%%%%%%%%%%%%%%%%%%%%%%%%%%
\subsection{Effect of High Utilization}
\label{subsec:high-utilization-effect}
%%%%%%%%%%%%%%%%%%%%%%%%%%%%%%%%%%%%%%%%%%%%%%%%%%%%%%%%%%%%

For \(M/M/1\),

\[
W
=
\frac1{\mu-\lambda}.
\]

As

\[
\lambda
\to
\mu^-,
\]

the denominator approaches zero.

Therefore,

\[
\boxed{
W\to\infty.
}
\]

Similarly,

\[
\boxed{
L\to\infty.
}
\]

This demonstrates why operating a variable service system near \(100\%\)
utilization can produce very long queues.

%%%%%%%%%%%%%%%%%%%%%%%%%%%%%%%%%%%%%%%%%%%%%%%%%%%%%%%%%%%%
\subsection{Worked Example: Utilization Sensitivity}
\label{subsec:worked-utilization-sensitivity}
%%%%%%%%%%%%%%%%%%%%%%%%%%%%%%%%%%%%%%%%%%%%%%%%%%%%%%%%%%%%

\begin{workedexamplebox}{Increasing Demand}

Suppose

\[
\mu=10.
\]

If

\[
\lambda=5,
\]

then

\[
W
=
\frac1{10-5}
=
0.2
\text{ hour}.
\]

If

\[
\lambda=9,
\]

then

\[
W
=
\frac1{10-9}
=
1
\text{ hour}.
\]

Thus increasing utilization from

\[
0.5
\]

to

\[
0.9
\]

increases mean system time from

\[
12
\]

minutes to

\[
60
\]

minutes.

\begin{answerbox}
\[
\boxed{
\text{Waiting grows nonlinearly as utilization approaches }1.
}
\]
\end{answerbox}

\end{workedexamplebox}

%%%%%%%%%%%%%%%%%%%%%%%%%%%%%%%%%%%%%%%%%%%%%%%%%%%%%%%%%%%%
\subsection{Distribution of \(M/M/1\) System Time}
\label{subsec:mm1-system-time-distribution}
%%%%%%%%%%%%%%%%%%%%%%%%%%%%%%%%%%%%%%%%%%%%%%%%%%%%%%%%%%%%

For a stable FCFS \(M/M/1\) queue, the total time \(W\) in the system has
an exponential distribution with rate

\[
\mu-\lambda.
\]

Thus,

\[
\boxed{
\mathbb{P}(W>t)
=
e^{-(\mu-\lambda)t}.
}
\]

Its mean is

\[
\boxed{
\frac1{\mu-\lambda}.
}
\]

%%%%%%%%%%%%%%%%%%%%%%%%%%%%%%%%%%%%%%%%%%%%%%%%%%%%%%%%%%%%
\subsection{\(M/M/1/K\) Finite-Capacity Queue}
\label{subsec:mm1k}
%%%%%%%%%%%%%%%%%%%%%%%%%%%%%%%%%%%%%%%%%%%%%%%%%%%%%%%%%%%%

In an \(M/M/1/K\) model, at most \(K\) customers may be in the system.

When

\[
N=K,
\]

new arrivals are blocked.

The birth rates are therefore

\[
\lambda_n
=
\begin{cases}
\lambda,
&
n<K,
\\
0,
&
n=K.
\end{cases}
\]

%%%%%%%%%%%%%%%%%%%%%%%%%%%%%%%%%%%%%%%%%%%%%%%%%%%%%%%%%%%%
\subsection{Stationary Distribution for \(M/M/1/K\)}
\label{subsec:mm1k-stationary}
%%%%%%%%%%%%%%%%%%%%%%%%%%%%%%%%%%%%%%%%%%%%%%%%%%%%%%%%%%%%

For

\[
\rho
=
\frac{\lambda}{\mu},
\]

the stationary probabilities satisfy

\[
p_n
=
\rho^n p_0.
\]

When

\[
\rho\neq1,
\]

normalization gives

\[
\boxed{
p_0
=
\frac{
1-\rho
}{
1-\rho^{K+1}
}.
}
\]

Hence,

\[
\boxed{
p_n
=
\frac{
(1-\rho)\rho^n
}{
1-\rho^{K+1}
}.
}
\]

%%%%%%%%%%%%%%%%%%%%%%%%%%%%%%%%%%%%%%%%%%%%%%%%%%%%%%%%%%%%
\subsection{Case \(\rho=1\) for \(M/M/1/K\)}
\label{subsec:mm1k-rho-one}
%%%%%%%%%%%%%%%%%%%%%%%%%%%%%%%%%%%%%%%%%%%%%%%%%%%%%%%%%%%%

If

\[
\rho=1,
\]

then all states have equal stationary probability:

\[
\boxed{
p_n
=
\frac1{K+1}.
}
\]

Because the system has finite capacity, a stationary distribution exists
even when

\[
\lambda\geq\mu.
\]

Blocked arrivals prevent unlimited growth.

%%%%%%%%%%%%%%%%%%%%%%%%%%%%%%%%%%%%%%%%%%%%%%%%%%%%%%%%%%%%
\subsection{Blocking Probability}
\label{subsec:blocking-probability}
%%%%%%%%%%%%%%%%%%%%%%%%%%%%%%%%%%%%%%%%%%%%%%%%%%%%%%%%%%%%

For \(M/M/1/K\), an arrival is blocked when the system is full.

By the Poisson-arrivals-see-time-averages property, abbreviated PASTA,
the blocking probability is

\[
\boxed{
P_{\mathrm{block}}
=
p_K.
}
\]

Thus the effective admitted arrival rate is

\[
\boxed{
\lambda_{\mathrm{eff}}
=
\lambda(1-p_K).
}
\]

%%%%%%%%%%%%%%%%%%%%%%%%%%%%%%%%%%%%%%%%%%%%%%%%%%%%%%%%%%%%
\subsection{Little's Law with Effective Throughput}
\label{subsec:littles-law-finite-capacity}
%%%%%%%%%%%%%%%%%%%%%%%%%%%%%%%%%%%%%%%%%%%%%%%%%%%%%%%%%%%%

In a finite-capacity system with blocked arrivals, Little's Law uses the
rate of customers actually entering the system:

\[
\boxed{
L
=
\lambda_{\mathrm{eff}}W.
}
\]

Using the offered arrival rate \(\lambda\) instead would overcount blocked
customers.

%%%%%%%%%%%%%%%%%%%%%%%%%%%%%%%%%%%%%%%%%%%%%%%%%%%%%%%%%%%%
\subsection{Multi-Server Queues}
\label{subsec:multi-server-queues}
%%%%%%%%%%%%%%%%%%%%%%%%%%%%%%%%%%%%%%%%%%%%%%%%%%%%%%%%%%%%

An \(M/M/c\) system has:

\[
\boxed{
\text{Poisson arrivals at rate }\lambda,
}
\]

\[
\boxed{
\text{exponential service rate }\mu
\text{ per server},
}
\]

and

\[
\boxed{
c
\text{ parallel servers}.
}
\]

%%%%%%%%%%%%%%%%%%%%%%%%%%%%%%%%%%%%%%%%%%%%%%%%%%%%%%%%%%%%
\subsection{Utilization in \(M/M/c\)}
\label{subsec:mmc-utilization}
%%%%%%%%%%%%%%%%%%%%%%%%%%%%%%%%%%%%%%%%%%%%%%%%%%%%%%%%%%%%

The total service capacity is

\[
c\mu.
\]

Define utilization per server as

\[
\boxed{
\rho
=
\frac{
\lambda
}{
c\mu
}.
}
\]

The stability condition is

\[
\boxed{
\rho<1,
}
\]

or equivalently,

\[
\boxed{
\lambda<c\mu.
}
\]

%%%%%%%%%%%%%%%%%%%%%%%%%%%%%%%%%%%%%%%%%%%%%%%%%%%%%%%%%%%%
\subsection{Offered Load}
\label{subsec:offered-load}
%%%%%%%%%%%%%%%%%%%%%%%%%%%%%%%%%%%%%%%%%%%%%%%%%%%%%%%%%%%%

Define the offered load

\[
\boxed{
a
=
\frac{\lambda}{\mu}.
}
\]

Then

\[
\boxed{
\rho
=
\frac{a}{c}.
}
\]

The quantity \(a\) is measured in Erlangs in telecommunications
applications.

%%%%%%%%%%%%%%%%%%%%%%%%%%%%%%%%%%%%%%%%%%%%%%%%%%%%%%%%%%%%
\subsection{Birth--Death Rates for \(M/M/c\)}
\label{subsec:mmc-birth-death}
%%%%%%%%%%%%%%%%%%%%%%%%%%%%%%%%%%%%%%%%%%%%%%%%%%%%%%%%%%%%

The birth rate is

\[
\boxed{
\lambda_n=\lambda.
}
\]

The service completion rate depends on how many servers are busy:

\[
\boxed{
\mu_n
=
\min\{n,c\}\mu.
}
\]

Thus:

\[
\mu_n
=
n\mu
\]

for

\[
n<c,
\]

and

\[
\mu_n
=
c\mu
\]

for

\[
n\geq c.
\]

%%%%%%%%%%%%%%%%%%%%%%%%%%%%%%%%%%%%%%%%%%%%%%%%%%%%%%%%%%%%
\subsection{Empty-System Probability for \(M/M/c\)}
\label{subsec:mmc-p0}
%%%%%%%%%%%%%%%%%%%%%%%%%%%%%%%%%%%%%%%%%%%%%%%%%%%%%%%%%%%%

Let

\[
a
=
\frac{\lambda}{\mu}.
\]

Then

\[
\boxed{
P_0
=
\left[
\sum_{n=0}^{c-1}
\frac{
a^n
}{
n!
}
+
\frac{
a^c
}{
c!(1-\rho)
}
\right]^{-1}.
}
\]

This normalizes the stationary distribution.

%%%%%%%%%%%%%%%%%%%%%%%%%%%%%%%%%%%%%%%%%%%%%%%%%%%%%%%%%%%%
\subsection{Probability of Waiting}
\label{subsec:erlang-c}
%%%%%%%%%%%%%%%%%%%%%%%%%%%%%%%%%%%%%%%%%%%%%%%%%%%%%%%%%%%%

In an \(M/M/c\) queue, an arriving customer waits if all \(c\) servers are
busy.

The Erlang C delay probability is

\[
\boxed{
P_{\mathrm{wait}}
=
\frac{
a^c
}{
c!(1-\rho)
}
P_0.
}
\]

This formula assumes an infinite waiting room and no abandonment.

%%%%%%%%%%%%%%%%%%%%%%%%%%%%%%%%%%%%%%%%%%%%%%%%%%%%%%%%%%%%
\subsection{Expected Queue Length for \(M/M/c\)}
\label{subsec:mmc-Lq}
%%%%%%%%%%%%%%%%%%%%%%%%%%%%%%%%%%%%%%%%%%%%%%%%%%%%%%%%%%%%

The mean number waiting is

\[
\boxed{
L_q
=
P_{\mathrm{wait}}
\frac{
\rho
}{
1-\rho
}.
}
\]

Equivalently,

\[
\boxed{
L_q
=
P_0
\frac{
a^c\rho
}{
c!(1-\rho)^2
}.
}
\]

%%%%%%%%%%%%%%%%%%%%%%%%%%%%%%%%%%%%%%%%%%%%%%%%%%%%%%%%%%%%
\subsection{Waiting Time for \(M/M/c\)}
\label{subsec:mmc-Wq}
%%%%%%%%%%%%%%%%%%%%%%%%%%%%%%%%%%%%%%%%%%%%%%%%%%%%%%%%%%%%

Using Little's Law,

\[
\boxed{
W_q
=
\frac{
L_q
}{
\lambda
}.
}
\]

The total expected time in the system is

\[
\boxed{
W
=
W_q
+
\frac1\mu.
}
\]

The expected number in the system is

\[
\boxed{
L
=
\lambda W.
}
\]

%%%%%%%%%%%%%%%%%%%%%%%%%%%%%%%%%%%%%%%%%%%%%%%%%%%%%%%%%%%%
\subsection{Worked Example: \(M/M/2\)}
\label{subsec:worked-mm2}
%%%%%%%%%%%%%%%%%%%%%%%%%%%%%%%%%%%%%%%%%%%%%%%%%%%%%%%%%%%%

\begin{workedexamplebox}{Two Parallel Servers}

Suppose

\[
\lambda=6
\]

customers per hour,

\[
\mu=4
\]

customers per hour per server, and

\[
c=2.
\]

Then

\[
a
=
\frac64
=
1.5,
\]

and

\[
\rho
=
\frac{1.5}{2}
=
0.75.
\]

First compute

\[
P_0
=
\left[
1
+
1.5
+
\frac{
1.5^2
}{
2(1-0.75)
}
\right]^{-1}.
\]

Since

\[
\frac{2.25}{0.5}
=
4.5,
\]

\[
P_0
=
\frac1{7}
\approx0.1429.
\]

The probability of waiting is

\[
P_{\mathrm{wait}}
=
4.5
\left(
\frac17
\right)
=
\frac{9}{14}
\approx0.6429.
\]

Then

\[
L_q
=
P_{\mathrm{wait}}
\frac{0.75}{0.25}
=
\frac{27}{14}
\approx1.9286.
\]

Thus,

\[
W_q
=
\frac{
1.9286
}{
6
}
\approx0.3214
\text{ hour}.
\]

Therefore,

\begin{answerbox}
\[
\boxed{
P_{\mathrm{wait}}
\approx0.6429,
}
\]

\[
\boxed{
W_q
\approx19.3
\text{ minutes}.
}
\]
\end{answerbox}

\end{workedexamplebox}

%%%%%%%%%%%%%%%%%%%%%%%%%%%%%%%%%%%%%%%%%%%%%%%%%%%%%%%%%%%%
\subsection{Erlang C Service-Level Formula}
\label{subsec:erlang-c-service-level}
%%%%%%%%%%%%%%%%%%%%%%%%%%%%%%%%%%%%%%%%%%%%%%%%%%%%%%%%%%%%

In an \(M/M/c\) queue, conditional on having to wait, queue waiting time is
exponential with rate

\[
c\mu-\lambda.
\]

Therefore,

\[
\boxed{
\mathbb{P}(W_q>t)
=
P_{\mathrm{wait}}
e^{-(c\mu-\lambda)t}.
}
\]

Hence,

\[
\boxed{
\mathbb{P}(W_q\leq t)
=
1
-
P_{\mathrm{wait}}
e^{-(c\mu-\lambda)t}.
}
\]

This formula is useful for staffing to service-level targets.

%%%%%%%%%%%%%%%%%%%%%%%%%%%%%%%%%%%%%%%%%%%%%%%%%%%%%%%%%%%%
\subsection{Loss Systems}
\label{subsec:loss-systems}
%%%%%%%%%%%%%%%%%%%%%%%%%%%%%%%%%%%%%%%%%%%%%%%%%%%%%%%%%%%%

In some systems, no waiting is allowed.

If all servers are busy, an arriving customer is immediately lost or
blocked.

A standard model is

\[
\boxed{
M/M/c/c.
}
\]

Examples include certain telephone and communication systems.

%%%%%%%%%%%%%%%%%%%%%%%%%%%%%%%%%%%%%%%%%%%%%%%%%%%%%%%%%%%%
\subsection{Erlang B Formula}
\label{subsec:erlang-b}
%%%%%%%%%%%%%%%%%%%%%%%%%%%%%%%%%%%%%%%%%%%%%%%%%%%%%%%%%%%%

Let

\[
a
=
\frac{\lambda}{\mu}.
\]

For an \(M/M/c/c\) loss system, the blocking probability is

\[
\boxed{
B(c,a)
=
\frac{
a^c/c!
}{
\displaystyle
\sum_{n=0}^{c}
a^n/n!
}.
}
\]

This is the Erlang B formula.

%%%%%%%%%%%%%%%%%%%%%%%%%%%%%%%%%%%%%%%%%%%%%%%%%%%%%%%%%%%%
\subsection{Erlang B Versus Erlang C}
\label{subsec:erlang-b-versus-c}
%%%%%%%%%%%%%%%%%%%%%%%%%%%%%%%%%%%%%%%%%%%%%%%%%%%%%%%%%%%%

Erlang B models:

\[
\boxed{
\text{no waiting}.
}
\]

Blocked arrivals leave immediately.

Erlang C models:

\[
\boxed{
\text{waiting allowed}.
}
\]

Customers wait when all servers are busy.

These models answer different operational questions.

%%%%%%%%%%%%%%%%%%%%%%%%%%%%%%%%%%%%%%%%%%%%%%%%%%%%%%%%%%%%
\subsection{\(M/G/1\) Queue}
\label{subsec:mg1-queue}
%%%%%%%%%%%%%%%%%%%%%%%%%%%%%%%%%%%%%%%%%%%%%%%%%%%%%%%%%%%%

An \(M/G/1\) queue has:

\[
\boxed{
\text{Poisson arrivals},
}
\]

\[
\boxed{
\text{general service-time distribution},
}
\]

and

\[
\boxed{
\text{one server}.
}
\]

Let

\[
\mathbb{E}[S]
=
\frac1\mu,
\]

and define

\[
\boxed{
\rho
=
\lambda
\mathbb{E}[S].
}
\]

Stability requires

\[
\boxed{
\rho<1.
}
\]

%%%%%%%%%%%%%%%%%%%%%%%%%%%%%%%%%%%%%%%%%%%%%%%%%%%%%%%%%%%%
\subsection{Pollaczek--Khinchine Mean Formula}
\label{subsec:pollaczek-khinchine}
%%%%%%%%%%%%%%%%%%%%%%%%%%%%%%%%%%%%%%%%%%%%%%%%%%%%%%%%%%%%

For a stable FCFS \(M/G/1\) queue,

\[
\boxed{
W_q
=
\frac{
\lambda
\mathbb{E}[S^2]
}{
2(1-\rho)
}.
}
\]

Then

\[
\boxed{
W
=
W_q
+
\mathbb{E}[S].
}
\]

Using Little's Law,

\[
\boxed{
L_q
=
\lambda W_q,
}
\]

and

\[
\boxed{
L
=
\lambda W.
}
\]

%%%%%%%%%%%%%%%%%%%%%%%%%%%%%%%%%%%%%%%%%%%%%%%%%%%%%%%%%%%%
\subsection{Role of Service Variability}
\label{subsec:service-variability}
%%%%%%%%%%%%%%%%%%%%%%%%%%%%%%%%%%%%%%%%%%%%%%%%%%%%%%%%%%%%

Since

\[
\mathbb{E}[S^2]
=
\operatorname{Var}(S)
+
(
\mathbb{E}[S]
)^2,
\]

greater service-time variance increases

\[
W_q.
\]

Thus two systems with the same average service time can have different
waiting times because of different variability.

%%%%%%%%%%%%%%%%%%%%%%%%%%%%%%%%%%%%%%%%%%%%%%%%%%%%%%%%%%%%
\subsection{Coefficient of Variation}
\label{subsec:service-coefficient-variation}
%%%%%%%%%%%%%%%%%%%%%%%%%%%%%%%%%%%%%%%%%%%%%%%%%%%%%%%%%%%%

Define the squared coefficient of variation of service time:

\[
\boxed{
c_s^2
=
\frac{
\operatorname{Var}(S)
}{
(
\mathbb{E}[S]
)^2
}.
}
\]

Then

\[
\boxed{
\mathbb{E}[S^2]
=
(1+c_s^2)
(
\mathbb{E}[S]
)^2.
}
\]

The Pollaczek--Khinchine formula becomes

\[
\boxed{
W_q
=
\frac{
\rho
}{
1-\rho
}
\frac{
1+c_s^2
}{
2
}
\mathbb{E}[S].
}
\]

%%%%%%%%%%%%%%%%%%%%%%%%%%%%%%%%%%%%%%%%%%%%%%%%%%%%%%%%%%%%
\subsection{\(M/D/1\) Queue}
\label{subsec:md1-queue}
%%%%%%%%%%%%%%%%%%%%%%%%%%%%%%%%%%%%%%%%%%%%%%%%%%%%%%%%%%%%

For deterministic service,

\[
\boxed{
S=\frac1\mu
}
\]

with zero variance.

Thus,

\[
\mathbb{E}[S^2]
=
\frac1{\mu^2}.
\]

The Pollaczek--Khinchine formula gives

\[
\boxed{
W_q
=
\frac{
\lambda
}{
2\mu^2(1-\rho)
}.
}
\]

Since

\[
\rho
=
\frac{\lambda}{\mu},
\]

\[
\boxed{
W_q
=
\frac{
\rho
}{
2\mu(1-\rho)
}.
}
\]

%%%%%%%%%%%%%%%%%%%%%%%%%%%%%%%%%%%%%%%%%%%%%%%%%%%%%%%%%%%%
\subsection{\(M/M/1\) as an \(M/G/1\) Special Case}
\label{subsec:mm1-as-mg1}
%%%%%%%%%%%%%%%%%%%%%%%%%%%%%%%%%%%%%%%%%%%%%%%%%%%%%%%%%%%%

For exponential service,

\[
\boxed{
\mathbb{E}[S^2]
=
\frac{2}{\mu^2}.
}
\]

Then

\[
W_q
=
\frac{
\lambda(2/\mu^2)
}{
2(1-\rho)
}.
\]

Therefore,

\[
\boxed{
W_q
=
\frac{
\lambda
}{
\mu^2(1-\rho)
}
=
\frac{
\lambda
}{
\mu(\mu-\lambda)
},
}
\]

which agrees with the \(M/M/1\) formula.

%%%%%%%%%%%%%%%%%%%%%%%%%%%%%%%%%%%%%%%%%%%%%%%%%%%%%%%%%%%%
\subsection{Worked Example: Effect of Service Variability}
\label{subsec:worked-service-variability}
%%%%%%%%%%%%%%%%%%%%%%%%%%%%%%%%%%%%%%%%%%%%%%%%%%%%%%%%%%%%

\begin{workedexamplebox}{Deterministic Versus Exponential Service}

Suppose

\[
\lambda=4
\]

per hour and

\[
\mu=6
\]

per hour.

Then

\[
\rho
=
\frac46
=
\frac23.
\]

For \(M/M/1\),

\[
W_q
=
\frac{
4
}{
6(6-4)
}
=
\frac13
\text{ hour}.
\]

Thus,

\[
W_q=20
\text{ minutes}.
\]

For \(M/D/1\),

\[
W_q
=
\frac{
\rho
}{
2\mu(1-\rho)
}.
\]

Hence,

\[
W_q
=
\frac{
2/3
}{
2(6)(1/3)
}
=
\frac16
\text{ hour}.
\]

Thus,

\[
W_q=10
\text{ minutes}.
\]

\begin{answerbox}
\[
\boxed{
\text{Removing service-time variability cuts the mean waiting time in half
in this example.}
}
\]
\end{answerbox}

\end{workedexamplebox}

%%%%%%%%%%%%%%%%%%%%%%%%%%%%%%%%%%%%%%%%%%%%%%%%%%%%%%%%%%%%
\subsection{\(G/G/1\) Queue}
\label{subsec:gg1-queue}
%%%%%%%%%%%%%%%%%%%%%%%%%%%%%%%%%%%%%%%%%%%%%%%%%%%%%%%%%%%%

A \(G/G/1\) queue allows:

\[
\boxed{
\text{general interarrival times},
}
\]

\[
\boxed{
\text{general service times},
}
\]

and

\[
\boxed{
\text{one server}.
}
\]

Exact formulas are usually much harder than for Markovian queues.

Approximation methods become important.

%%%%%%%%%%%%%%%%%%%%%%%%%%%%%%%%%%%%%%%%%%%%%%%%%%%%%%%%%%%%
\subsection{Arrival Variability}
\label{subsec:arrival-variability}
%%%%%%%%%%%%%%%%%%%%%%%%%%%%%%%%%%%%%%%%%%%%%%%%%%%%%%%%%%%%

Let \(A\) be the interarrival time.

Define

\[
\boxed{
c_a^2
=
\frac{
\operatorname{Var}(A)
}{
(
\mathbb{E}[A]
)^2
}.
}
\]

For exponential arrivals,

\[
\boxed{
c_a^2=1.
}
\]

For deterministic arrivals,

\[
\boxed{
c_a^2=0.
}
\]

%%%%%%%%%%%%%%%%%%%%%%%%%%%%%%%%%%%%%%%%%%%%%%%%%%%%%%%%%%%%
\subsection{Kingman's Approximation}
\label{subsec:kingman-approximation}
%%%%%%%%%%%%%%%%%%%%%%%%%%%%%%%%%%%%%%%%%%%%%%%%%%%%%%%%%%%%

A widely used approximation for a \(G/G/1\) FCFS queue is

\[
\boxed{
W_q
\approx
\frac{
\rho
}{
1-\rho
}
\cdot
\frac{
c_a^2+c_s^2
}{
2
}
\cdot
\mathbb{E}[S].
}
\]

This is often called Kingman's approximation or the VUT approximation.

%%%%%%%%%%%%%%%%%%%%%%%%%%%%%%%%%%%%%%%%%%%%%%%%%%%%%%%%%%%%
\subsection{Variability--Utilization--Time Interpretation}
\label{subsec:vut-interpretation}
%%%%%%%%%%%%%%%%%%%%%%%%%%%%%%%%%%%%%%%%%%%%%%%%%%%%%%%%%%%%

Kingman's approximation can be viewed as

\[
\boxed{
\text{waiting}
\approx
\text{utilization factor}
\times
\text{variability factor}
\times
\text{service-time scale}.
}
\]

Specifically,

\[
\boxed{
\frac{\rho}{1-\rho}
}
\]

captures utilization,

\[
\boxed{
\frac{c_a^2+c_s^2}{2}
}
\]

captures variability, and

\[
\boxed{
\mathbb{E}[S]
}
\]

sets the time scale.

%%%%%%%%%%%%%%%%%%%%%%%%%%%%%%%%%%%%%%%%%%%%%%%%%%%%%%%%%%%%
\subsection{Operational Meaning of Kingman's Approximation}
\label{subsec:kingman-operational-meaning}
%%%%%%%%%%%%%%%%%%%%%%%%%%%%%%%%%%%%%%%%%%%%%%%%%%%%%%%%%%%%

Waiting can be reduced by:

\[
\boxed{
\text{reducing utilization},
}
\]

\[
\boxed{
\text{reducing arrival variability},
}
\]

or

\[
\boxed{
\text{reducing service variability}.
}
\]

This gives practical guidance beyond simply increasing the average service
rate.

%%%%%%%%%%%%%%%%%%%%%%%%%%%%%%%%%%%%%%%%%%%%%%%%%%%%%%%%%%%%
\subsection{Finite-Source Queues}
\label{subsec:finite-source-queues}
%%%%%%%%%%%%%%%%%%%%%%%%%%%%%%%%%%%%%%%%%%%%%%%%%%%%%%%%%%%%

Suppose only \(N\) machines can fail and request repair.

If \(n\) machines are already failed, only

\[
N-n
\]

machines remain capable of generating new repair requests.

The arrival rate therefore depends on system state.

This differs from infinite-source Poisson queues.

%%%%%%%%%%%%%%%%%%%%%%%%%%%%%%%%%%%%%%%%%%%%%%%%%%%%%%%%%%%%
\subsection{Machine-Repair Model Preview}
\label{subsec:machine-repair-model}
%%%%%%%%%%%%%%%%%%%%%%%%%%%%%%%%%%%%%%%%%%%%%%%%%%%%%%%%%%%%

Suppose each operating machine fails independently at rate \(\lambda\).

If \(n\) machines are currently failed, then the total failure rate is

\[
\boxed{
\lambda_n
=
(N-n)\lambda.
}
\]

If one repair server works at rate \(\mu\),

\[
\boxed{
\mu_n
=
\mu,
\qquad
n\geq1.
}
\]

This is another birth--death process.

%%%%%%%%%%%%%%%%%%%%%%%%%%%%%%%%%%%%%%%%%%%%%%%%%%%%%%%%%%%%
\subsection{Transient Versus Steady-State Analysis}
\label{subsec:transient-steady-state-queueing}
%%%%%%%%%%%%%%%%%%%%%%%%%%%%%%%%%%%%%%%%%%%%%%%%%%%%%%%%%%%%

Steady-state analysis studies long-run behavior such as

\[
\boxed{
\lim_{t\to\infty}
\mathbb{P}
(
N(t)=n
).
}
\]

Transient analysis studies behavior at finite time \(t\).

A system may take substantial time to approach equilibrium.

For short operating periods, transient analysis can be more relevant.

%%%%%%%%%%%%%%%%%%%%%%%%%%%%%%%%%%%%%%%%%%%%%%%%%%%%%%%%%%%%
\subsection{PASTA Property}
\label{subsec:pasta-property}
%%%%%%%%%%%%%%%%%%%%%%%%%%%%%%%%%%%%%%%%%%%%%%%%%%%%%%%%%%%%

PASTA stands for

\[
\boxed{
\text{Poisson Arrivals See Time Averages}.
}
\]

For a system with Poisson arrivals, the probability an arriving customer
observes the system in a particular state equals the long-run fraction of
time the system spends in that state, under suitable standard conditions.

This simplifies many arrival-based performance calculations.

%%%%%%%%%%%%%%%%%%%%%%%%%%%%%%%%%%%%%%%%%%%%%%%%%%%%%%%%%%%%
\subsection{Balking}
\label{subsec:balking}
%%%%%%%%%%%%%%%%%%%%%%%%%%%%%%%%%%%%%%%%%%%%%%%%%%%%%%%%%%%%

Balking occurs when an arriving customer decides not to enter the queue.

The decision may depend on:

\[
\boxed{
\text{queue length},
}
\]

\[
\boxed{
\text{expected delay},
}
\]

or

\[
\boxed{
\text{customer preference}.
}
\]

Balking makes the effective arrival process state dependent.

%%%%%%%%%%%%%%%%%%%%%%%%%%%%%%%%%%%%%%%%%%%%%%%%%%%%%%%%%%%%
\subsection{Reneging}
\label{subsec:reneging}
%%%%%%%%%%%%%%%%%%%%%%%%%%%%%%%%%%%%%%%%%%%%%%%%%%%%%%%%%%%%

Reneging occurs when a customer joins the queue but leaves before receiving
service.

The time a customer is willing to wait is called patience time.

Abandonment can significantly affect call centers and healthcare queues.

%%%%%%%%%%%%%%%%%%%%%%%%%%%%%%%%%%%%%%%%%%%%%%%%%%%%%%%%%%%%
\subsection{Jockeying}
\label{subsec:jockeying}
%%%%%%%%%%%%%%%%%%%%%%%%%%%%%%%%%%%%%%%%%%%%%%%%%%%%%%%%%%%%

Jockeying occurs when a customer switches from one queue to another.

This is relevant when parallel servers maintain separate waiting lines.

Jockeying creates dependence between queues and complicates analysis.

%%%%%%%%%%%%%%%%%%%%%%%%%%%%%%%%%%%%%%%%%%%%%%%%%%%%%%%%%%%%
\subsection{Single Queue Versus Separate Queues}
\label{subsec:single-versus-separate-queues}
%%%%%%%%%%%%%%%%%%%%%%%%%%%%%%%%%%%%%%%%%%%%%%%%%%%%%%%%%%%%

With several servers, customers may wait in:

\[
\boxed{
\text{one common queue},
}
\]

or

\[
\boxed{
\text{separate queues}.
}
\]

A common queue often pools variability and improves load balancing.

Separate queues may be operationally simpler but can leave one server idle
while another queue remains long.

%%%%%%%%%%%%%%%%%%%%%%%%%%%%%%%%%%%%%%%%%%%%%%%%%%%%%%%%%%%%
\subsection{Pooling Effect}
\label{subsec:queue-pooling-effect}
%%%%%%%%%%%%%%%%%%%%%%%%%%%%%%%%%%%%%%%%%%%%%%%%%%%%%%%%%%%%

Pooling customers into one queue can reduce inefficiency caused by random
differences among individual queues.

The principle is:

\[
\boxed{
\text{shared capacity can absorb variability more effectively}.
}
\]

This is a recurring idea in service-system design.

%%%%%%%%%%%%%%%%%%%%%%%%%%%%%%%%%%%%%%%%%%%%%%%%%%%%%%%%%%%%
\subsection{Priority Queues}
\label{subsec:priority-queues}
%%%%%%%%%%%%%%%%%%%%%%%%%%%%%%%%%%%%%%%%%%%%%%%%%%%%%%%%%%%%

A priority queue divides customers into classes.

Higher-priority customers are selected before lower-priority customers.

Applications include:

\[
\boxed{
\text{emergency medicine},
}
\]

\[
\boxed{
\text{computer scheduling},
}
\]

and

\[
\boxed{
\text{communication traffic}.
}
\]

%%%%%%%%%%%%%%%%%%%%%%%%%%%%%%%%%%%%%%%%%%%%%%%%%%%%%%%%%%%%
\subsection{Nonpreemptive Priority}
\label{subsec:nonpreemptive-priority}
%%%%%%%%%%%%%%%%%%%%%%%%%%%%%%%%%%%%%%%%%%%%%%%%%%%%%%%%%%%%

Under nonpreemptive priority, a high-priority arrival cannot interrupt a
customer already in service.

It only moves ahead of lower-priority customers waiting in the queue.

%%%%%%%%%%%%%%%%%%%%%%%%%%%%%%%%%%%%%%%%%%%%%%%%%%%%%%%%%%%%
\subsection{Preemptive Priority}
\label{subsec:preemptive-priority}
%%%%%%%%%%%%%%%%%%%%%%%%%%%%%%%%%%%%%%%%%%%%%%%%%%%%%%%%%%%%

Under preemptive priority, a sufficiently high-priority arrival can
interrupt current service.

The displaced customer may later resume or restart service depending on
the model.

Preemption can greatly reduce delay for high-priority classes.

%%%%%%%%%%%%%%%%%%%%%%%%%%%%%%%%%%%%%%%%%%%%%%%%%%%%%%%%%%%%
\subsection{Conservation of Work}
\label{subsec:queue-conservation-work}
%%%%%%%%%%%%%%%%%%%%%%%%%%%%%%%%%%%%%%%%%%%%%%%%%%%%%%%%%%%%

Priority rules redistribute waiting time among customer classes.

They do not create service capacity.

Reducing waiting for one class generally comes at the expense of another
class unless additional resources or other efficiencies are introduced.

%%%%%%%%%%%%%%%%%%%%%%%%%%%%%%%%%%%%%%%%%%%%%%%%%%%%%%%%%%%%
\subsection{Queueing Networks}
\label{subsec:queueing-networks}
%%%%%%%%%%%%%%%%%%%%%%%%%%%%%%%%%%%%%%%%%%%%%%%%%%%%%%%%%%%%

A queueing network contains several service stations through which
customers move.

After completing service at station \(i\), a customer may:

\[
\boxed{
\text{leave the system},
}
\]

or

\[
\boxed{
\text{move to another station}.
}
\]

Applications include production lines, communication networks, and
computer systems.

%%%%%%%%%%%%%%%%%%%%%%%%%%%%%%%%%%%%%%%%%%%%%%%%%%%%%%%%%%%%
\subsection{Open Queueing Networks}
\label{subsec:open-queueing-networks}
%%%%%%%%%%%%%%%%%%%%%%%%%%%%%%%%%%%%%%%%%%%%%%%%%%%%%%%%%%%%

An open queueing network allows customers to enter from and leave to the
external environment.

Let

\[
\gamma_i
\]

be the external arrival rate to station \(i\).

Let

\[
p_{ij}
\]

be the probability that a customer completing service at station \(i\)
moves next to station \(j\).

%%%%%%%%%%%%%%%%%%%%%%%%%%%%%%%%%%%%%%%%%%%%%%%%%%%%%%%%%%%%
\subsection{Traffic Equations}
\label{subsec:queue-network-traffic-equations}
%%%%%%%%%%%%%%%%%%%%%%%%%%%%%%%%%%%%%%%%%%%%%%%%%%%%%%%%%%%%

Let

\[
\lambda_j
\]

be the total arrival rate to station \(j\).

Then the traffic equations are

\[
\boxed{
\lambda_j
=
\gamma_j
+
\sum_i
\lambda_i p_{ij}.
}
\]

In vector form,

\[
\boxed{
\boldsymbol{\lambda}
=
\boldsymbol{\gamma}
+
P^T
\boldsymbol{\lambda}.
}
\]

Thus,

\[
\boxed{
(I-P^T)
\boldsymbol{\lambda}
=
\boldsymbol{\gamma}.
}
\]

%%%%%%%%%%%%%%%%%%%%%%%%%%%%%%%%%%%%%%%%%%%%%%%%%%%%%%%%%%%%
\subsection{Jackson Networks}
\label{subsec:jackson-networks}
%%%%%%%%%%%%%%%%%%%%%%%%%%%%%%%%%%%%%%%%%%%%%%%%%%%%%%%%%%%%

An open Jackson network has standard assumptions including:

\[
\boxed{
\text{external Poisson arrivals},
}
\]

\[
\boxed{
\text{exponential service},
}
\]

and

\[
\boxed{
\text{Markovian routing}.
}
\]

Under stability conditions, the joint stationary distribution has a
product form.

This makes otherwise complicated networks analytically tractable.

%%%%%%%%%%%%%%%%%%%%%%%%%%%%%%%%%%%%%%%%%%%%%%%%%%%%%%%%%%%%
\subsection{Jackson Network Stability}
\label{subsec:jackson-network-stability}
%%%%%%%%%%%%%%%%%%%%%%%%%%%%%%%%%%%%%%%%%%%%%%%%%%%%%%%%%%%%

If station \(i\) is a single exponential server with service rate
\(\mu_i\), then the station utilization is

\[
\boxed{
\rho_i
=
\frac{
\lambda_i
}{
\mu_i
}.
}
\]

A standard stability requirement is

\[
\boxed{
\rho_i<1
}
\]

for every station.

The most highly utilized station may act as a bottleneck.

%%%%%%%%%%%%%%%%%%%%%%%%%%%%%%%%%%%%%%%%%%%%%%%%%%%%%%%%%%%%
\subsection{Worked Example: Network Traffic Equations}
\label{subsec:worked-network-traffic}
%%%%%%%%%%%%%%%%%%%%%%%%%%%%%%%%%%%%%%%%%%%%%%%%%%%%%%%%%%%%

\begin{workedexamplebox}{Two-Station Queueing Network}

Suppose external arrivals enter station \(1\) at rate

\[
\gamma_1=4,
\]

and station \(2\) at rate

\[
\gamma_2=1.
\]

After station \(1\), \(25\%\) of customers go to station \(2\).

After station \(2\), \(20\%\) return to station \(1\).

Thus,

\[
\lambda_1
=
4
+
0.2\lambda_2,
\]

and

\[
\lambda_2
=
1
+
0.25\lambda_1.
\]

Substitute:

\[
\lambda_1
=
4
+
0.2
\left(
1+0.25\lambda_1
\right).
\]

Then

\[
\lambda_1
=
4.2
+
0.05\lambda_1.
\]

Therefore,

\[
0.95\lambda_1
=
4.2,
\]

so

\[
\lambda_1
=
\frac{84}{19}
\approx4.421.
\]

Then

\[
\lambda_2
=
1
+
0.25
\left(
\frac{84}{19}
\right)
=
\frac{40}{19}
\approx2.105.
\]

\begin{answerbox}
\[
\boxed{
\lambda_1
=
\frac{84}{19}
\approx4.421,
}
\]

\[
\boxed{
\lambda_2
=
\frac{40}{19}
\approx2.105.
}
\]
\end{answerbox}

\end{workedexamplebox}

%%%%%%%%%%%%%%%%%%%%%%%%%%%%%%%%%%%%%%%%%%%%%%%%%%%%%%%%%%%%
\subsection{Closed Queueing Networks}
\label{subsec:closed-queueing-networks}
%%%%%%%%%%%%%%%%%%%%%%%%%%%%%%%%%%%%%%%%%%%%%%%%%%%%%%%%%%%%

A closed queueing network has a fixed population circulating among service
stations.

No customers enter or leave externally.

Examples include:

\[
\boxed{
\text{a fixed fleet of vehicles},
}
\]

or

\[
\boxed{
\text{a fixed number of jobs circulating among machines}.
}
\]

Closed networks often require specialized product-form or numerical
methods.

%%%%%%%%%%%%%%%%%%%%%%%%%%%%%%%%%%%%%%%%%%%%%%%%%%%%%%%%%%%%
\subsection{Bottlenecks in Queueing Networks}
\label{subsec:queueing-network-bottlenecks}
%%%%%%%%%%%%%%%%%%%%%%%%%%%%%%%%%%%%%%%%%%%%%%%%%%%%%%%%%%%%

A bottleneck station has comparatively high workload relative to capacity.

As system demand increases, its utilization reaches \(1\) first.

The bottleneck can dominate:

\[
\boxed{
\text{throughput},
}
\]

\[
\boxed{
\text{queue length},
}
\]

and

\[
\boxed{
\text{delay}.
}
\]

%%%%%%%%%%%%%%%%%%%%%%%%%%%%%%%%%%%%%%%%%%%%%%%%%%%%%%%%%%%%
\subsection{Service-Level Measures}
\label{subsec:queue-service-level-measures}
%%%%%%%%%%%%%%%%%%%%%%%%%%%%%%%%%%%%%%%%%%%%%%%%%%%%%%%%%%%%

Operational performance may be evaluated by:

\[
\boxed{
\mathbb{E}[W_q],
}
\]

\[
\boxed{
\mathbb{P}(W_q\leq t),
}
\]

\[
\boxed{
\mathbb{P}(N_q\leq k),
}
\]

\[
\boxed{
\text{blocking probability},
}
\]

or

\[
\boxed{
\text{abandonment probability}.
}
\]

The appropriate metric depends on the application.

%%%%%%%%%%%%%%%%%%%%%%%%%%%%%%%%%%%%%%%%%%%%%%%%%%%%%%%%%%%%
\subsection{Staffing Decisions}
\label{subsec:queue-staffing}
%%%%%%%%%%%%%%%%%%%%%%%%%%%%%%%%%%%%%%%%%%%%%%%%%%%%%%%%%%%%

Increasing the number of servers reduces delay but increases staffing cost.

A basic staffing objective might balance

\[
\boxed{
\text{server cost}
}
\]

and

\[
\boxed{
\text{waiting cost}.
}
\]

For \(c\) servers, a simple model is

\[
\boxed{
C(c)
=
cC_s
+
C_wL_q(c),
}
\]

where

\[
C_s
=
\text{cost per server},
\]

and

\[
C_w
=
\text{waiting cost per customer-time unit}.
\]

%%%%%%%%%%%%%%%%%%%%%%%%%%%%%%%%%%%%%%%%%%%%%%%%%%%%%%%%%%%%
\subsection{Worked Example: Staffing Tradeoff}
\label{subsec:worked-staffing-tradeoff}
%%%%%%%%%%%%%%%%%%%%%%%%%%%%%%%%%%%%%%%%%%%%%%%%%%%%%%%%%%%%

\begin{workedexamplebox}{Cost of Waiting Versus Service Capacity}

Suppose adding another server costs

\[
\$30
\]

per hour.

Suppose the additional server reduces the average number waiting from

\[
4.2
\]

to

\[
1.1.
\]

If each waiting customer costs

\[
\$12
\]

per customer-hour, the waiting-cost reduction is

\[
12(4.2-1.1)
=
12(3.1)
=
37.2.
\]

Since

\[
37.2>30,
\]

the extra server reduces the modeled total hourly cost by

\[
37.2-30
=
7.2.
\]

\begin{answerbox}
\[
\boxed{
\text{Under this cost model, adding the server is beneficial.}
}
\]
\end{answerbox}

\end{workedexamplebox}

%%%%%%%%%%%%%%%%%%%%%%%%%%%%%%%%%%%%%%%%%%%%%%%%%%%%%%%%%%%%
\subsection{Economies of Scale in Service Capacity}
\label{subsec:service-economies-scale}
%%%%%%%%%%%%%%%%%%%%%%%%%%%%%%%%%%%%%%%%%%%%%%%%%%%%%%%%%%%%

Pooling demand across shared servers can improve utilization while
maintaining service levels.

A larger pooled system may simultaneously achieve:

\[
\boxed{
\text{higher utilization}
}
\]

and

\[
\boxed{
\text{lower probability of extreme delay}
}
\]

than several isolated small systems.

This is an important service-system design principle.

%%%%%%%%%%%%%%%%%%%%%%%%%%%%%%%%%%%%%%%%%%%%%%%%%%%%%%%%%%%%
\subsection{Queueing and Scheduling}
\label{subsec:queueing-scheduling-connection}
%%%%%%%%%%%%%%%%%%%%%%%%%%%%%%%%%%%%%%%%%%%%%%%%%%%%%%%%%%%%

Queueing theory treats arrival and service uncertainty.

Scheduling theory often assumes a set of known jobs and asks how they
should be ordered or assigned.

Thus:

\[
\boxed{
\text{queueing}
\rightarrow
\text{stochastic congestion},
}
\]

while

\[
\boxed{
\text{scheduling}
\rightarrow
\text{planned sequencing and allocation}.
}
\]

The two areas overlap in stochastic scheduling problems.

%%%%%%%%%%%%%%%%%%%%%%%%%%%%%%%%%%%%%%%%%%%%%%%%%%%%%%%%%%%%
\subsection{Queueing and Inventory}
\label{subsec:queueing-inventory-connection}
%%%%%%%%%%%%%%%%%%%%%%%%%%%%%%%%%%%%%%%%%%%%%%%%%%%%%%%%%%%%

Both queueing and inventory involve flow through systems.

In queueing:

\[
\boxed{
\text{excess arrivals}
\rightarrow
\text{waiting customers}.
}
\]

In inventory:

\[
\boxed{
\text{excess demand}
\rightarrow
\text{shortages or backorders}.
}
\]

Both study the interaction among rates, variability, and capacity.

%%%%%%%%%%%%%%%%%%%%%%%%%%%%%%%%%%%%%%%%%%%%%%%%%%%%%%%%%%%%
\subsection{Simulation of Queues}
\label{subsec:queue-simulation}
%%%%%%%%%%%%%%%%%%%%%%%%%%%%%%%%%%%%%%%%%%%%%%%%%%%%%%%%%%%%

When analytical formulas are unavailable, discrete-event simulation can
estimate queue performance.

A simulation tracks events such as:

\[
\boxed{
\text{arrival},
}
\]

\[
\boxed{
\text{service start},
}
\]

and

\[
\boxed{
\text{service completion}.
}
\]

Repeated simulation estimates waiting times, queue lengths, and service
levels.

%%%%%%%%%%%%%%%%%%%%%%%%%%%%%%%%%%%%%%%%%%%%%%%%%%%%%%%%%%%%
\subsection{Why Simulation Is Useful}
\label{subsec:why-queue-simulation}
%%%%%%%%%%%%%%%%%%%%%%%%%%%%%%%%%%%%%%%%%%%%%%%%%%%%%%%%%%%%

Simulation can handle:

\[
\boxed{
\text{nonexponential service times},
}
\]

\[
\boxed{
\text{time-varying arrival rates},
}
\]

\[
\boxed{
\text{complex priority rules},
}
\]

\[
\boxed{
\text{abandonment},
}
\]

and

\[
\boxed{
\text{multiple interacting resources}.
}
\]

Its disadvantage is that simulation estimates performance rather than
providing exact symbolic formulas.

%%%%%%%%%%%%%%%%%%%%%%%%%%%%%%%%%%%%%%%%%%%%%%%%%%%%%%%%%%%%
\subsection{Warm-Up and Initial Conditions}
\label{subsec:queue-simulation-warmup}
%%%%%%%%%%%%%%%%%%%%%%%%%%%%%%%%%%%%%%%%%%%%%%%%%%%%%%%%%%%%

A simulation started with an empty system may initially behave differently
from a steady-state queue.

This initial transient can bias estimates.

A warm-up period may be discarded before steady-state statistics are
collected.

%%%%%%%%%%%%%%%%%%%%%%%%%%%%%%%%%%%%%%%%%%%%%%%%%%%%%%%%%%%%
\subsection{Time-Varying Arrival Rates}
\label{subsec:time-varying-arrival-rates}
%%%%%%%%%%%%%%%%%%%%%%%%%%%%%%%%%%%%%%%%%%%%%%%%%%%%%%%%%%%%

Real systems often have arrival rates

\[
\boxed{
\lambda(t)
}
\]

that vary with time.

Examples include:

\[
\boxed{
\text{morning and evening transportation peaks},
}
\]

or

\[
\boxed{
\text{hourly call-center demand}.
}
\]

Stationary \(M/M/c\) formulas may be inadequate when demand changes
rapidly.

%%%%%%%%%%%%%%%%%%%%%%%%%%%%%%%%%%%%%%%%%%%%%%%%%%%%%%%%%%%%
\subsection{Nonhomogeneous Poisson Process Preview}
\label{subsec:nonhomogeneous-poisson-queue}
%%%%%%%%%%%%%%%%%%%%%%%%%%%%%%%%%%%%%%%%%%%%%%%%%%%%%%%%%%%%

A nonhomogeneous Poisson process has time-dependent rate

\[
\lambda(t).
\]

The expected number of arrivals over \([a,b]\) is

\[
\boxed{
\int_a^b
\lambda(t)\,dt.
}
\]

Such arrival models are useful for time-dependent service systems.

%%%%%%%%%%%%%%%%%%%%%%%%%%%%%%%%%%%%%%%%%%%%%%%%%%%%%%%%%%%%
\subsection{Queueing Optimization}
\label{subsec:queueing-optimization}
%%%%%%%%%%%%%%%%%%%%%%%%%%%%%%%%%%%%%%%%%%%%%%%%%%%%%%%%%%%%

Queueing models can be embedded inside optimization problems.

Decision variables may include:

\[
\boxed{
\text{number of servers},
}
\]

\[
\boxed{
\text{service rates},
}
\]

\[
\boxed{
\text{capacity},
}
\]

\[
\boxed{
\text{routing},
}
\]

or

\[
\boxed{
\text{priority rules}.
}
\]

The objective may balance cost and service performance.

%%%%%%%%%%%%%%%%%%%%%%%%%%%%%%%%%%%%%%%%%%%%%%%%%%%%%%%%%%%%
\subsection{Capacity Planning}
\label{subsec:queue-capacity-planning}
%%%%%%%%%%%%%%%%%%%%%%%%%%%%%%%%%%%%%%%%%%%%%%%%%%%%%%%%%%%%

Suppose system demand is expected to increase.

Queueing models can estimate how much additional service capacity is
required to maintain a target such as

\[
\boxed{
\mathbb{E}[W_q]
\leq
w_{\max},
}
\]

or

\[
\boxed{
\mathbb{P}(W_q\leq t_0)
\geq
\alpha.
}
\]

This converts queueing analysis into a design problem.

%%%%%%%%%%%%%%%%%%%%%%%%%%%%%%%%%%%%%%%%%%%%%%%%%%%%%%%%%%%%
\subsection{Common Errors}
\label{subsec:queueing-common-errors}
%%%%%%%%%%%%%%%%%%%%%%%%%%%%%%%%%%%%%%%%%%%%%%%%%%%%%%%%%%%%

\subsubsection{Using Arrival Rate and Mean Interarrival Time Interchangeably}

If

\[
\lambda
=
6
\text{ customers per hour},
\]

then mean interarrival time is

\[
\boxed{
\frac16
\text{ hour}
=
10
\text{ minutes}.
}
\]

The units must remain consistent.

%%%%%%%%%%%%%%%%%%%%%%%%%%%%%%%%%%%%%%%%%%%%%%%%%%%%%%%%%%%%
\subsubsection{Using Service Time Instead of Service Rate}

If mean service time is \(5\) minutes, then the service rate is

\[
\boxed{
\mu
=
12
\text{ customers per hour}.
}
\]

The value \(5\) is not the service rate.

%%%%%%%%%%%%%%%%%%%%%%%%%%%%%%%%%%%%%%%%%%%%%%%%%%%%%%%%%%%%
\subsubsection{Forgetting the Stability Condition}

The infinite-capacity \(M/M/1\) steady-state formulas require

\[
\boxed{
\lambda<\mu.
}
\]

Using them when \(\lambda\geq\mu\) is invalid.

%%%%%%%%%%%%%%%%%%%%%%%%%%%%%%%%%%%%%%%%%%%%%%%%%%%%%%%%%%%%
\subsubsection{Using \(\rho=\lambda/\mu\) for a Multi-Server Queue}

For \(M/M/c\),

\[
\boxed{
\rho
=
\frac{\lambda}{c\mu}.
}
\]

The total service capacity is \(c\mu\).

%%%%%%%%%%%%%%%%%%%%%%%%%%%%%%%%%%%%%%%%%%%%%%%%%%%%%%%%%%%%
\subsubsection{Confusing \(L\) and \(L_q\)}

\[
L
=
\text{mean number in the entire system},
\]

while

\[
L_q
=
\text{mean number waiting}.
\]

Customers currently receiving service are included in \(L\) but not
\(L_q\).

%%%%%%%%%%%%%%%%%%%%%%%%%%%%%%%%%%%%%%%%%%%%%%%%%%%%%%%%%%%%
\subsubsection{Confusing \(W\) and \(W_q\)}

\[
W
=
\text{mean total time in system},
\]

while

\[
W_q
=
\text{mean waiting time before service}.
\]

For a single service stage,

\[
\boxed{
W
=
W_q+\mathbb{E}[S].
}
\]

%%%%%%%%%%%%%%%%%%%%%%%%%%%%%%%%%%%%%%%%%%%%%%%%%%%%%%%%%%%%
\subsubsection{Using Offered Arrivals in Little's Law When Customers Are Blocked}

For finite-capacity systems, use the effective admitted throughput:

\[
\boxed{
L
=
\lambda_{\mathrm{eff}}W.
}
\]

Blocked arrivals never enter the system.

%%%%%%%%%%%%%%%%%%%%%%%%%%%%%%%%%%%%%%%%%%%%%%%%%%%%%%%%%%%%
\subsubsection{Assuming Little's Law Requires an \(M/M/1\) Queue}

Little's Law is much more general.

It does not require Poisson arrivals or exponential service.

%%%%%%%%%%%%%%%%%%%%%%%%%%%%%%%%%%%%%%%%%%%%%%%%%%%%%%%%%%%%
\subsubsection{Ignoring Variability}

Mean arrival and service rates alone may not determine waiting performance.

Service-time and interarrival-time variability can have major effects.

%%%%%%%%%%%%%%%%%%%%%%%%%%%%%%%%%%%%%%%%%%%%%%%%%%%%%%%%%%%%
\subsubsection{Confusing Erlang B and Erlang C}

Erlang B assumes blocked customers are lost.

Erlang C assumes customers wait.

The formulas are not interchangeable.

%%%%%%%%%%%%%%%%%%%%%%%%%%%%%%%%%%%%%%%%%%%%%%%%%%%%%%%%%%%%
\subsubsection{Assuming Poisson Arrivals in Every Real System}

Poisson models can be useful approximations, but real arrival processes may
have:

\[
\boxed{
\text{time dependence},
}
\]

\[
\boxed{
\text{burstiness},
}
\]

or

\[
\boxed{
\text{correlation}.
}
\]

Model assumptions should be checked against data.

%%%%%%%%%%%%%%%%%%%%%%%%%%%%%%%%%%%%%%%%%%%%%%%%%%%%%%%%%%%%
\subsubsection{Ignoring Customer Behavior}

Balking, abandonment, priorities, and jockeying can materially change
performance.

A simple \(M/M/c\) formula may then be inadequate.

%%%%%%%%%%%%%%%%%%%%%%%%%%%%%%%%%%%%%%%%%%%%%%%%%%%%%%%%%%%%
\subsubsection{Treating Steady-State Results as Immediate}

Steady-state formulas describe long-run behavior.

A newly opened or short-duration system may still be in a transient phase.

%%%%%%%%%%%%%%%%%%%%%%%%%%%%%%%%%%%%%%%%%%%%%%%%%%%%%%%%%%%%
\subsubsection{Assuming Maximum Utilization Is Operationally Optimal}

Higher utilization improves resource use but usually increases waiting
nonlinearly.

Capacity planning must balance both effects.

%%%%%%%%%%%%%%%%%%%%%%%%%%%%%%%%%%%%%%%%%%%%%%%%%%%%%%%%%%%%
\subsection{Queueing-Theory Workflow}
\label{subsec:queueing-workflow}
%%%%%%%%%%%%%%%%%%%%%%%%%%%%%%%%%%%%%%%%%%%%%%%%%%%%%%%%%%%%

A practical workflow is:

\begin{enumerate}
    \item identify the customer or job class;
    \item identify arrival times or arrival rates;
    \item identify service-time distributions;
    \item count the number of servers;
    \item identify queue capacity;
    \item identify the service discipline;
    \item determine whether arrivals are approximately Poisson;
    \item determine whether service times are approximately exponential;
    \item choose appropriate Kendall notation;
    \item compute utilization;
    \item check stability;
    \item identify whether exact formulas are available;
    \item compute \(L\), \(L_q\), \(W\), and \(W_q\);
    \item use Little's Law to cross-check results;
    \item evaluate blocking or service-level probabilities when relevant;
    \item examine variability;
    \item identify bottlenecks;
    \item consider simulation for complex systems;
    \item compare alternative staffing or capacity choices;
    \item interpret results in operational units.
\end{enumerate}

%%%%%%%%%%%%%%%%%%%%%%%%%%%%%%%%%%%%%%%%%%%%%%%%%%%%%%%%%%%%
\subsection{Choosing a Queueing Model}
\label{subsec:choosing-queueing-model}
%%%%%%%%%%%%%%%%%%%%%%%%%%%%%%%%%%%%%%%%%%%%%%%%%%%%%%%%%%%%

For one exponential server with Poisson arrivals, use:

\[
\boxed{
M/M/1.
}
\]

For several parallel exponential servers, use:

\[
\boxed{
M/M/c.
}
\]

For a finite waiting room, consider:

\[
\boxed{
M/M/1/K
}
\]

or the corresponding multi-server finite-capacity model.

For no waiting, consider:

\[
\boxed{
M/M/c/c.
}
\]

For general service times with Poisson arrivals and one server, consider:

\[
\boxed{
M/G/1.
}
\]

For general arrivals and service, consider:

\[
\boxed{
G/G/1
}
\]

approximations or simulation.

%%%%%%%%%%%%%%%%%%%%%%%%%%%%%%%%%%%%%%%%%%%%%%%%%%%%%%%%%%%%
\subsection{Applications}
\label{subsec:queueing-applications}
%%%%%%%%%%%%%%%%%%%%%%%%%%%%%%%%%%%%%%%%%%%%%%%%%%%%%%%%%%%%

\begin{applicationbox}

\textbf{Healthcare.}
Patients wait for emergency rooms, diagnostic equipment, physicians, and
beds.

\medskip

\textbf{Airports and Transportation.}
Aircraft, passengers, vehicles, and cargo may queue for gates, security,
runways, loading equipment, and service facilities.

\medskip

\textbf{Manufacturing.}
Jobs wait for machines, inspection stations, repair facilities, and
assembly resources.

\medskip

\textbf{Call Centers.}
Customers wait for available agents while managers choose staffing levels
to satisfy service targets.

\medskip

\textbf{Computing.}
Tasks, requests, and packets wait for processors, servers, storage devices,
and communication channels.

\medskip

\textbf{Telecommunications.}
Queueing models describe packet congestion, channel blocking, and call
capacity.

\medskip

\textbf{Retail and Service Systems.}
Checkout lanes, restaurants, banks, and customer-service desks require
capacity decisions under random demand.

\medskip

\textbf{Supply Chains.}
Congestion can arise at loading docks, warehouses, inspection facilities,
and transportation terminals.

\end{applicationbox}

%%%%%%%%%%%%%%%%%%%%%%%%%%%%%%%%%%%%%%%%%%%%%%%%%%%%%%%%%%%%
\subsection{Exercises}
\label{subsec:queueing-exercises}
%%%%%%%%%%%%%%%%%%%%%%%%%%%%%%%%%%%%%%%%%%%%%%%%%%%%%%%%%%%%

\begin{exercise}[1]
Define a queueing system.
\end{exercise}

\begin{exercise}[1]
Define arrival rate.
\end{exercise}

\begin{exercise}[1]
Define service rate.
\end{exercise}

\begin{exercise}[1]
Define utilization.
\end{exercise}

\begin{exercise}[1]
Interpret \(M/M/1\).
\end{exercise}

\begin{exercise}[1]
Interpret \(M/M/c\).
\end{exercise}

\begin{exercise}[1]
State Little's Law.
\end{exercise}

\begin{exercise}[1]
Define balking.
\end{exercise}

\begin{exercise}[1]
Define reneging.
\end{exercise}

\begin{exercise}[1]
Distinguish Erlang B from Erlang C.
\end{exercise}

\begin{exercise}[2]
If customers arrive every \(8\) minutes on average, find the arrival rate
per hour.
\end{exercise}

\begin{exercise}[2]
If mean service time is \(12\) minutes, find the service rate per hour.
\end{exercise}

\begin{exercise}[2]
For

\[
\lambda=4,
\qquad
\mu=7,
\]

compute \(\rho\).
\end{exercise}

\begin{exercise}[2]
Determine whether the previous \(M/M/1\) queue is stable.
\end{exercise}

\begin{exercise}[2]
For \(\rho=0.75\), compute \(P_0\).
\end{exercise}

\begin{exercise}[2]
For \(\rho=0.6\), compute

\[
P(N=3).
\]
\end{exercise}

\begin{exercise}[2]
For \(\rho=0.8\), compute

\[
P(N\geq4).
\]
\end{exercise}

\begin{exercise}[3]
For an \(M/M/1\) queue with

\[
\lambda=5,
\qquad
\mu=8,
\]

compute

\[
L,
\quad
L_q,
\quad
W,
\quad
W_q.
\]
\end{exercise}

\begin{exercise}[3]
Verify Little's Law for the previous problem.
\end{exercise}

\begin{exercise}[3]
Convert the waiting-time answers in the previous problem from hours to
minutes.
\end{exercise}

\begin{exercise}[3]
Compare \(W\) for

\[
\lambda=4,\mu=10
\]

and

\[
\lambda=9,\mu=10.
\]
\end{exercise}

\begin{exercise}[3]
Explain why waiting time rises sharply as utilization approaches one.
\end{exercise}

\begin{exercise}[3]
Derive

\[
P_n
=
(1-\rho)\rho^n
\]

for \(M/M/1\).
\end{exercise}

\begin{exercise}[3]
Derive

\[
L
=
\frac{\rho}{1-\rho}.
\]
\end{exercise}

\begin{exercise}[3]
Use Little's Law to derive \(W\) from \(L\).
\end{exercise}

\begin{exercise}[3]
Use

\[
W=W_q+\frac1\mu
\]

to verify the \(M/M/1\) waiting-time formulas.
\end{exercise}

\begin{exercise}[3]
For an \(M/M/1/5\) queue with specified \(\lambda\) and \(\mu\), compute
the stationary probabilities.
\end{exercise}

\begin{exercise}[3]
Compute the blocking probability for the previous finite-capacity queue.
\end{exercise}

\begin{exercise}[3]
Compute the effective arrival rate for the same queue.
\end{exercise}

\begin{exercise}[3]
Explain why finite capacity allows a steady state even when
\(\lambda\geq\mu\).
\end{exercise}

\begin{exercise}[3]
For an \(M/M/3\) queue, compute

\[
\rho
=
\frac{\lambda}{3\mu}.
\]
\end{exercise}

\begin{exercise}[3]
Compute \(P_0\) for a small \(M/M/2\) system.
\end{exercise}

\begin{exercise}[3]
Compute the Erlang C probability of waiting for the previous system.
\end{exercise}

\begin{exercise}[3]
Use Little's Law to compute \(W_q\) from \(L_q\).
\end{exercise}

\begin{exercise}[3]
Compute a service-level probability

\[
P(W_q\leq t)
\]

for an \(M/M/c\) model.
\end{exercise}

\begin{exercise}[3]
Compute an Erlang B blocking probability for a small loss system.
\end{exercise}

\begin{exercise}[3]
Explain why Erlang B does not include queue waiting time.
\end{exercise}

\begin{exercise}[3]
Use the Pollaczek--Khinchine formula for an \(M/G/1\) queue.
\end{exercise}

\begin{exercise}[3]
Compare waiting times for exponential and deterministic service with equal
means.
\end{exercise}

\begin{exercise}[3]
Compute \(c_s^2\) from a given service-time variance and mean.
\end{exercise}

\begin{exercise}[3]
Use Kingman's approximation for a \(G/G/1\) queue.
\end{exercise}

\begin{exercise}[3]
Explain how arrival variability affects waiting time.
\end{exercise}

\begin{exercise}[3]
Explain how service variability affects waiting time.
\end{exercise}

\begin{exercise}[3]
Construct the birth and death rates for a finite-source machine-repair
model.
\end{exercise}

\begin{exercise}[3]
Solve a two-station queueing-network traffic equation.
\end{exercise}

\begin{exercise}[3]
Identify the bottleneck station from given arrival and service rates.
\end{exercise}

\begin{exercise}[3]
Formulate a simple staffing-cost model using \(L_q(c)\).
\end{exercise}

\begin{exercise}[3]
Compare one pooled \(M/M/2\) queue with two separate \(M/M/1\) queues
conceptually.
\end{exercise}

\begin{exercise}[3]
Give examples of balking, reneging, and jockeying.
\end{exercise}

\begin{exercise}[3]
Explain the difference between preemptive and nonpreemptive priority.
\end{exercise}

\begin{exercise}[3]
Explain why a steady-state formula may be inappropriate for a short
operating period.
\end{exercise}

\begin{exercise}[3]
Design a basic discrete-event simulation for an \(M/M/1\) queue.
\end{exercise}

\begin{exercise}[4]
Derive the stationary distribution of a general birth--death process.
\end{exercise}

\begin{exercise}[4]
Prove the \(M/M/1\) steady-state formulas.
\end{exercise}

\begin{exercise}[4]
Derive the distribution of \(M/M/1\) system time.
\end{exercise}

\begin{exercise}[4]
Study PASTA formally.
\end{exercise}

\begin{exercise}[4]
Derive the stationary distribution of \(M/M/1/K\).
\end{exercise}

\begin{exercise}[4]
Derive the Erlang C formula for \(M/M/c\).
\end{exercise}

\begin{exercise}[4]
Derive the Erlang B loss formula.
\end{exercise}

\begin{exercise}[4]
Study recursive evaluation of Erlang B.
\end{exercise}

\begin{exercise}[4]
Derive the Pollaczek--Khinchine transform formula for \(M/G/1\).
\end{exercise}

\begin{exercise}[4]
Derive the Pollaczek--Khinchine mean formula.
\end{exercise}

\begin{exercise}[4]
Study residual service time in \(M/G/1\).
\end{exercise}

\begin{exercise}[4]
Study the inspection paradox in queueing systems.
\end{exercise}

\begin{exercise}[4]
Derive Kingman's heavy-traffic approximation.
\end{exercise}

\begin{exercise}[4]
Study \(GI/M/1\) queues.
\end{exercise}

\begin{exercise}[4]
Study priority \(M/G/1\) queues.
\end{exercise}

\begin{exercise}[4]
Study processor-sharing queues.
\end{exercise}

\begin{exercise}[4]
Study queues with abandonment.
\end{exercise}

\begin{exercise}[4]
Study \(M/M/c+M\) models for call centers.
\end{exercise}

\begin{exercise}[4]
Derive Jackson-network product-form probabilities.
\end{exercise}

\begin{exercise}[4]
Study closed Gordon--Newell networks.
\end{exercise}

\begin{exercise}[4]
Study mean-value analysis for closed queueing networks.
\end{exercise}

\begin{exercise}[4]
Investigate heavy-traffic limits and diffusion approximations.
\end{exercise}

\begin{exercise}[4]
Study fluid approximations for large queueing systems.
\end{exercise}

\begin{exercise}[4]
Investigate queueing-control and dynamic staffing problems.
\end{exercise}

%%%%%%%%%%%%%%%%%%%%%%%%%%%%%%%%%%%%%%%%%%%%%%%%%%%%%%%%%%%%
\subsection{Conceptual Summary}
\label{subsec:queueing-summary}
%%%%%%%%%%%%%%%%%%%%%%%%%%%%%%%%%%%%%%%%%%%%%%%%%%%%%%%%%%%%

Queueing theory studies stochastic congestion in service systems.

For an \(M/M/1\) queue,

\[
\boxed{
\rho
=
\frac{\lambda}{\mu}
<
1.
}
\]

The steady-state distribution is

\[
\boxed{
P_n
=
(1-\rho)\rho^n.
}
\]

The principal performance measures are

\[
\boxed{
L
=
\frac{\rho}{1-\rho},
}
\]

\[
\boxed{
L_q
=
\frac{\rho^2}{1-\rho},
}
\]

\[
\boxed{
W
=
\frac1{\mu-\lambda},
}
\]

and

\[
\boxed{
W_q
=
\frac{\lambda}{
\mu(\mu-\lambda)
}.
}
\]

Little's Law gives

\[
\boxed{
L=\lambda W,
}
\]

and

\[
\boxed{
L_q=\lambda W_q.
}
\]

For multiple servers,

\[
\boxed{
\rho
=
\frac{\lambda}{c\mu}.
}
\]

Erlang C describes waiting systems, while Erlang B describes loss systems.

For \(M/G/1\),

\[
\boxed{
W_q
=
\frac{
\lambda\mathbb{E}[S^2]
}{
2(1-\rho)
}.
}
\]

Thus queueing performance depends not only on average rates, but also on

\[
\boxed{
\text{utilization}
+
\text{variability}
+
\text{capacity}.
}
\]

Queueing theory therefore combines

\[
\boxed{
\text{probability}
+
\text{stochastic processes}
+
\text{operations research}
+
\text{service-system design}.
}
\]

%%%%%%%%%%%%%%%%%%%%%%%%%%%%%%%%%%%%%%%%%%%%%%%%%%%%%%%%%%%%
\subsection{Section Connections}
\label{subsec:queueing-connections}
%%%%%%%%%%%%%%%%%%%%%%%%%%%%%%%%%%%%%%%%%%%%%%%%%%%%%%%%%%%%

\chapterconnection{
Queueing Theory applies probability and stochastic-process models to
congestion, waiting, and service capacity. The birth--death structure of
Markovian queues connects directly with Chapter~7, while staffing and
capacity decisions connect queueing analysis with Optimization and
Stochastic Optimization. Queueing networks extend the network ideas of
Section~\ref{sec:network-optimization} from deterministic flow to random
arrivals and service. The next section, Scheduling Theory, studies how a
known collection of jobs should be assigned and sequenced across resources
to optimize completion time, lateness, flow time, and related performance
measures.
}

%%%%%%%%%%%%%%%%%%%%%%%%%%%%%%%%%%%%%%%%%%%%%%%%%%%%%%%%%%%%
% SECTION SOURCE TRACKING
%%%%%%%%%%%%%%%%%%%%%%%%%%%%%%%%%%%%%%%%%%%%%%%%%%%%%%%%%%%%

%------------------------------------------------------------
% 10.11 Queueing Theory
%------------------------------------------------------------
%
% Source basis:
%   - Standard queueing theory.
%   - Standard birth--death and Markovian queue models.
%   - Standard operations-research service-system models.
%
% Used for:
%   - queueing-system definitions;
%   - arrivals and service processes;
%   - arrival and service rates;
%   - queue capacities;
%   - calling populations;
%   - queue disciplines;
%   - Kendall notation;
%   - Poisson arrivals;
%   - exponential interarrival and service times;
%   - memoryless property;
%   - utilization and traffic intensity;
%   - birth--death models;
%   - M/M/1 stationary distribution;
%   - queue-length probabilities;
%   - L, L_q, W, and W_q;
%   - Little's Law;
%   - stability;
%   - finite-capacity M/M/1/K queues;
%   - blocking and effective arrival rates;
%   - M/M/c systems;
%   - Erlang C;
%   - service-level calculations;
%   - Erlang B loss systems;
%   - M/G/1 queues;
%   - Pollaczek--Khinchine formulas;
%   - service variability;
%   - M/D/1;
%   - G/G/1 and Kingman approximation;
%   - finite-source models;
%   - transient and steady-state analysis;
%   - PASTA;
%   - balking, reneging, and jockeying;
%   - pooling;
%   - priority queues;
%   - queueing networks;
%   - traffic equations;
%   - Jackson networks;
%   - bottlenecks;
%   - staffing;
%   - simulation;
%   - time-varying arrivals;
%   - queueing optimization;
%   - applications and exercises.
%
% Notes:
%   - Scheduling Theory is developed next in Section 10.12.
%   - Inventory Theory follows in Section 10.13.
%   - Markov-chain and stochastic-process foundations are developed in
%     Chapter 7.
%
%%%%%%%%%%%%%%%%%%%%%%%%%%%%%%%%%%%%%%%%%%%%%%%%%%%%%%%%%%%%